% LEGACY, UNCOMPILED DRAFT FRAGMENT. The authoritative manuscript is ../main.tex.
\begin{abstract}
The BraTS local-synthesis task asks for plausible tumor-free tissue to be generated
inside a masked region of a T1-weighted brain MRI while preserving the observed
anatomy. We present CATCH, a conditional 3D diffusion framework operating on an
invertible Haar-wavelet representation. Its denoiser receives noisy target
subbands, voided-image subbands, and a signed mask. A tumor-excluded wavelet loss
with a hole-localized term provides supervision, while hard compositing preserves
every observed voxel. We compare fixed-mask training with strict random
tumor-component augmentation and a configured weighted mixture of tumor-derived,
blob, and ellipsoid masks. Arm-specific checkpoints and trajectory policies are
selected on development data, then frozen for a held-out 75-case confirmation
restricted to artificially removed healthy tissue. The selected weighted-mixture
pipeline (five trajectories, voxel-wise mean) achieved SSIM
$0.7972\pm0.1308$, PSNR $19.18\pm1.80$\,dB, and MSE
$0.009874\pm0.005359$ (mean $\pm$ sample SD). Against compute-matched random
augmentation, its mean improvements were $0.0188$ SSIM (95\% bootstrap CI:
$0.0127$--$0.0249$), $0.948$\,dB PSNR, and $0.002641$ MSE reduction; all three
corrected paired tests favored weighted mixture. Fixed-mask single-trajectory
inference achieved SSIM $0.7561\pm0.1615$. These within-CATCH results support the
complete weighted-mixture policy, but one training seed, arm-specific checkpoint
selection, and unequal trajectory counts relative to fixed mask limit causal and
cross-method conclusions.

\keywords{MRI inpainting \and Brain MRI \and Medical image synthesis
\and Diffusion models \and Wavelet transform \and BraTS}
\end{abstract}
